%% kapa: exam=1 round=36 year=2025 period=1 form=A subject=civil_law subject_ko=민법 seq=10 number=10
%% source: Q-Net 감정평가사 기출문제 (cst00309) artlSeq=5231525 file=raw/1008/5231525/2025년 제36회 감정평가사 1차시험 1교시 문제지 원본.pdf
%% conversion: kordoc 4.24.51 -> tools/gen1.py (auto); review=proofread-sonnet; flags=none
\begin{question}{10}
사기ㆍ강박에 의한 의사표시에 관한 설명으로 옳은 것을 모두 고른 것은? (다툼이 있으면 판례에 따름)
\begin{framedtext}
ㄱ. 상대방의 대리인에 의한 강박은 제3자의 강박에 해당한다. \par
ㄴ. 제3자가 사기에 의한 의사표시의 취소에 대항하기 위해서는 특별한 사정이 없는 한 자신의 선의를 증명해야 한다. \par
ㄷ. 거래의 중요한 사항에 관한 사실을 신의성실의 의무에 비추어 비난받을 정도의 방법으로 허위로 고지한 것은 기망행위에 해당한다.
\end{framedtext}
\begin{choices}
  \item ㄱ
  \item ㄷ
  \item ㄱ, ㄴ
  \item ㄴ, ㄷ
  \item ㄱ, ㄴ, ㄷ
\end{choices}
\end{question}
\answer{2}
